\documentclass[11pt,a4paper]{article}
\usepackage{amsmath}
\usepackage{amssymb}
\usepackage{geometry}
\usepackage{listings}
\usepackage{xcolor}
\usepackage{hyperref}

\geometry{margin=1in}

\lstset{
    language=Python,
    basicstyle=\ttfamily\small,
    keywordstyle=\color{blue},
    commentstyle=\color{gray},
    stringstyle=\color{red},
    breaklines=true,
    showstringspaces=false,
    backgroundcolor=\color{lightgray!20}
}

\title{Comprehensive Deep Learning Notes: 3Blue1Brown Series}
\author{Based on Neural Networks Chapters 1-3}
\date{}

\begin{document}

\maketitle

\section{Video 1: Neural Networks Structure (Chapter 1)}

\subsection{Overview}
Introduction to handwritten digit recognition using neural networks. This chapter covers the basic architecture: neurons, layers, weights, and biases, along with the mathematical foundation of forward propagation.

\subsection{Key Concepts}

\subsubsection{Neurons and Activation}
\begin{itemize}
    \item \textbf{Neuron}: A computational unit that holds a number between 0 and 1
    \item \textbf{Activation}: The value inside a neuron, representing how ``lit up'' it is
    \item Input layer: 784 neurons (28×28 pixel grid)
    \item Output layer: 10 neurons (digits 0-9)
    \item Hidden layers: Arbitrary choice (2 layers with 16 neurons each in example)
\end{itemize}

\subsubsection{Network Parameters}
\begin{itemize}
    \item Total neurons: $784 \to 16 \to 16 \to 10$
    \item Total weights and biases: $\approx 13,000$ parameters
\end{itemize}

\subsection{Mathematical Foundation}

\subsubsection{Forward Propagation}
For a neuron in layer $l+1$:

\textbf{Weighted Sum Computation:}
\[
z_j^{(l+1)} = \sum_{i=1}^{n_l} w_{ij}^{(l)} a_i^{(l)} + b_j^{(l)}
\]

Where:
\begin{itemize}
    \item $w_{ij}^{(l)}$: weight from neuron $i$ in layer $l$ to neuron $j$ in layer $l+1$
    \item $a_i^{(l)}$: activation of neuron $i$ in layer $l$
    \item $b_j^{(l)}$: bias for neuron $j$ in layer $l$
    \item $z_j^{(l+1)}$: weighted sum before activation function
\end{itemize}

\textbf{Activation Function (Sigmoid):}
\[
a_j^{(l+1)} = \sigma(z_j^{(l+1)}) = \frac{1}{1 + e^{-z_j^{(l+1)}}}
\]

Alternatively, ReLU (modern approach):
\[
a_j^{(l+1)} = \max(0, z_j^{(l+1)})
\]

\subsubsection{Matrix Form}
\[
\mathbf{a}^{(l+1)} = \sigma\left(\mathbf{W}^{(l)} \mathbf{a}^{(l)} + \mathbf{b}^{(l)}\right)
\]

Where:
\begin{itemize}
    \item $\mathbf{W}^{(l)}$: weight matrix (rows = connections to one neuron in next layer)
    \item $\mathbf{a}^{(l)}$: activation vector from previous layer
    \item $\mathbf{b}^{(l)}$: bias vector
    \item $\sigma$: activation function applied element-wise
\end{itemize}

\subsection{Why Layered Structure?}
\begin{enumerate}
    \item Layer $1 \to 2$: Detect edges and small patterns
    \item Layer $2 \to 3$: Combine edges to recognize loops and lines
    \item Layer $3 \to$ Output: Recognize complete digits from patterns
\end{enumerate}

\subsection{Code Example}
\begin{lstlisting}
import numpy as np

def sigmoid(z):
    return 1 / (1 + np.exp(-z))

def forward_pass(X, weights, biases):
    """Forward pass through network"""
    a = X
    for W, b in zip(weights, biases):
        z = np.dot(W, a) + b
        a = sigmoid(z)
    return a

# Network initialization
input_size = 784
hidden_sizes = [16, 16]
output_size = 10

weights = [np.random.randn(hidden_sizes[0], input_size) * 0.01]
for i in range(len(hidden_sizes)-1):
    weights.append(np.random.randn(
        hidden_sizes[i+1], hidden_sizes[i]) * 0.01)
weights.append(np.random.randn(output_size, hidden_sizes[-1]) * 0.01)

biases = [np.zeros((size,)) for size in hidden_sizes + [output_size]]

# Forward pass
output = forward_pass(image, weights, biases)
predicted_digit = np.argmax(output)
\end{lstlisting}

\newpage

\section{Video 2: Gradient Descent and Learning (Chapter 2)}

\subsection{Overview}
This chapter introduces the cost function definition, gradient descent algorithm, and how networks learn through optimization.

\subsection{Cost Function}

\subsubsection{Single Training Example}
\[
C_{\text{single}} = \sum_{k=1}^{10} (y_k - a_k^{(L)})^2
\]

Where:
\begin{itemize}
    \item $y_k$: desired output for digit $k$ (1 if correct digit, 0 otherwise)
    \item $a_k^{(L)}$: actual output activation
    \item $L$: final layer
    \item Sum over all 10 output neurons
\end{itemize}

\subsubsection{Total Cost Function}
\[
C_{\text{total}} = \frac{1}{n} \sum_{i=1}^{n} C_i
\]

Where:
\begin{itemize}
    \item $n$: number of training examples
    \item $C_i$: cost for $i$-th training example
\end{itemize}

\subsection{Gradient Descent}

\subsubsection{Conceptual Framework}
\begin{enumerate}
    \item Start with random weights and biases
    \item Compute gradient of cost function: $\nabla C$
    \item Update parameters in direction of negative gradient
    \item Repeat until convergence
\end{enumerate}

\subsubsection{Gradient Vector}
\[
\nabla C = \left[\frac{\partial C}{\partial w_1}, \frac{\partial C}{\partial w_2}, \ldots, \frac{\partial C}{\partial w_{13000}}\right]
\]

Each component tells us:
\begin{itemize}
    \item \textbf{Sign}: Whether to increase or decrease parameter
    \item \textbf{Magnitude}: How much that parameter affects cost (relative importance)
\end{itemize}

\subsubsection{Update Rule}
\[
w_{\text{new}} = w_{\text{old}} - \alpha \frac{\partial C}{\partial w}
\]
\[
b_{\text{new}} = b_{\text{old}} - \alpha \frac{\partial C}{\partial b}
\]

Where:
\begin{itemize}
    \item $\alpha$: learning rate (small positive number, e.g., 0.01)
    \item $\frac{\partial C}{\partial w}$, $\frac{\partial C}{\partial b}$: partial derivatives (components of gradient)
\end{itemize}

\subsection{Multivariable Calculus Connection}
For 2D example:
\[
\nabla C = \left[\frac{\partial C}{\partial x}, \frac{\partial C}{\partial y}\right]
\]

\begin{itemize}
    \item Points in direction of steepest ascent
    \item Negative gradient points downhill (minimization direction)
    \item Magnitude indicates steepness of slope
\end{itemize}

\subsection{Code Example}
\begin{lstlisting}
def compute_cost(output, target):
    """Compute MSE cost"""
    return np.sum((target - output)**2)

def gradient_descent_step(weights, biases, X, y, 
                         learning_rate):
    """Single gradient descent step"""
    # Forward pass
    activations = forward_pass(X[0], weights, biases)
    
    # Compute cost
    cost = compute_cost(activations, y[0])
    
    # Note: Gradient computation requires backpropagation
    # (see Chapter 3)
    
    return weights, biases, cost

# Training loop
for epoch in range(num_epochs):
    for batch_X, batch_y in mini_batches:
        weights, biases, cost = gradient_descent_step(
            weights, biases, batch_X, batch_y, 
            learning_rate=0.01
        )
    print(f"Epoch {epoch}, Cost: {cost}")
\end{lstlisting}

\subsection{Key Insights}
\begin{itemize}
    \item Network learns $\approx 96\%$ accuracy on test data with simple 2-hidden-layer architecture
    \item Can improve to 98\% with tweaking
    \item \textbf{Important caveat}: Network doesn't necessarily learn patterns we expect
    \begin{itemize}
        \item Second layer weights appear almost random
        \item Network can confidently classify random noise
        \item Cost function doesn't incentivize semantic features
    \end{itemize}
\end{itemize}

\subsection{Stochastic Gradient Descent (SGD)}

Problems with full gradient descent:
\begin{itemize}
    \item Computing gradient over ALL training examples is slow
    \item $10,000\text{s of examples} \times 13,000\text{ parameters} = \text{very expensive}$
\end{itemize}

Solution:
\begin{enumerate}
    \item Randomly shuffle training data
    \item Divide into mini-batches (e.g., 100 examples each)
    \item Compute gradient using mini-batch only
    \item Update weights
    \item Repeat with next mini-batch
\end{enumerate}

Effect:
\begin{itemize}
    \item ``Drunk man stumbling downhill'' vs ``carefully calculated descent''
    \item Quick steps with approximate gradients vs slow steps with exact gradient
    \item Usually converges faster in practice
\end{itemize}

\newpage

\section{Video 3: Backpropagation (Chapter 3)}

\subsection{Overview}
This chapter covers how to efficiently compute gradients using the backpropagation algorithm, the chain rule application in neural networks, and visual intuition before mathematical formalism.

\subsection{Key Concept: What Should Change?}

\subsubsection{Output Layer Perspective}
For each output neuron, we want:
\begin{itemize}
    \item Desired neuron: activation $\to 1$
    \item Other neurons: activation $\to 0$
\end{itemize}

Desired changes (nudges):
\begin{itemize}
    \item If $\text{output}_k < \text{target}_k$: nudge UP
    \item If $\text{output}_k > \text{target}_k$: nudge DOWN
    \item Magnitude of nudge $\propto |\text{output}_k - \text{target}_k|$
\end{itemize}

\subsubsection{Three Ways to Achieve Nudge}
For neuron $j$ with activation $a_j$:
\[
z_j = \sum_i w_{ij} a_i + b_j
\]
\[
a_j = \sigma(z_j)
\]

To increase $a_j$:
\begin{enumerate}
    \item \textbf{Increase bias}: $\Delta b_j > 0$
    \item \textbf{Increase weights}: $\Delta w_{ij} > 0$ (especially if $a_i$ is large)
    \item \textbf{Change previous activations}: $\Delta a_i > 0$ (for positive weights) or $< 0$ (for negative weights)
\end{enumerate}

\subsection{Backpropagation Algorithm (Intuitive)}

\subsubsection{Step 1: Output Layer}
\[
\text{error}_{\text{output}} = (a_{\text{output}} - y_{\text{target}})
\]
\[
\text{desired\_change}_{\text{output}} = -2 \cdot \text{error}_{\text{output}}
\]

\subsubsection{Step 2: Propagate Backward}
For each neuron in previous layer:
\[
\text{desire}_i = \sum_j w_{ij} \cdot \text{desire}_j
\]

Key insight: ``Neurons that fire together wire together''
\begin{itemize}
    \item Strong connections (large $w_{ij}$) propagate more desire
    \item Active neurons (large $a_i$) cause larger weight updates
\end{itemize}

\subsubsection{Step 3: Update Weights and Biases}
\[
\Delta w_{ij} \propto \text{desire}_j \times a_i
\]

Intuition:
\begin{itemize}
    \item If neuron $j$ needs to change a lot (large desire)
    \item AND neuron $i$ is very active (large $a_i$)
    \item THEN strengthen connection $w_{ij}$
\end{itemize}

\subsection{Backpropagation Algorithm (Mathematical)}

\subsubsection{Chain Rule Application}
For cost $C$ and weight $w_{ij}$:
\[
\frac{\partial C}{\partial w_{ij}} = \frac{\partial C}{\partial a_j} \times \frac{\partial a_j}{\partial z_j} \times \frac{\partial z_j}{\partial w_{ij}}
\]

Where:
\begin{itemize}
    \item $\frac{\partial z_j}{\partial w_{ij}} = a_i$
    \item For sigmoid: $\frac{\partial a_j}{\partial z_j} = \sigma'(z_j) = \sigma(z_j)(1 - \sigma(z_j))$
    \item For ReLU: $\frac{\partial a_j}{\partial z_j} = 1$ if $z_j > 0$, else $0$
\end{itemize}

Define: $\delta_j = \frac{\partial C}{\partial z_j}$ (error in $z$)

Then:
\[
\frac{\partial C}{\partial w_{ij}} = \delta_j \times a_i
\]
\[
\frac{\partial C}{\partial b_j} = \delta_j
\]

\subsubsection{Recursive Definition of $\delta$}
For output layer $L$:
\[
\delta_j^{(L)} = (a_j^{(L)} - y_j) \times \sigma'(z_j^{(L)})
\]

For hidden layer $l$:
\[
\delta_j^{(l)} = \left(\sum_k w_{jk}^{(l)} \delta_k^{(l+1)}\right) \times \sigma'(z_j^{(l)})
\]

Note: $\delta_k^{(l+1)}$ comes from next layer. This enables efficient backward pass.

\subsubsection{Full Backpropagation Update}

\begin{enumerate}
    \item Forward pass: compute all $a^{(l)}$ and $z^{(l)}$

    \item Compute output layer error:
    \[
    \boldsymbol{\delta}^{(L)} = (a^{(L)} - \mathbf{y}) \odot \sigma'(\mathbf{z}^{(L)})
    \]
    where $\odot$ is element-wise multiplication

    \item Backward pass ($l$ from $L-1$ down to $1$):
    \[
    \boldsymbol{\delta}^{(l)} = (\mathbf{W}^{(l)T} \boldsymbol{\delta}^{(l+1)}) \odot \sigma'(\mathbf{z}^{(l)})
    \]

    \item Compute gradients:
    \[
    \frac{\partial C}{\partial \mathbf{W}^{(l)}} = \boldsymbol{\delta}^{(l+1)} (\mathbf{a}^{(l)})^T
    \]
    \[
    \frac{\partial C}{\partial \mathbf{b}^{(l)}} = \boldsymbol{\delta}^{(l+1)}
    \]

    \item Update parameters:
    \[
    \mathbf{W}^{(l)} := \mathbf{W}^{(l)} - \alpha \frac{\partial C}{\partial \mathbf{W}^{(l)}}
    \]
    \[
    \mathbf{b}^{(l)} := \mathbf{b}^{(l)} - \alpha \frac{\partial C}{\partial \mathbf{b}^{(l)}}
    \]
\end{enumerate}

\subsection{Code Example}
\begin{lstlisting}
def sigmoid_derivative(a):
    """Derivative of sigmoid: sigma'(z) = sigma(z)(1-sigma(z))"""
    return a * (1 - a)

def relu_derivative(z):
    """Derivative of ReLU"""
    return (z > 0).astype(float)

def backpropagation(X, y, weights, biases, activations, 
                    z_values):
    """
    X: input (784,)
    y: target (10,)
    weights, biases: network parameters
    activations: list of a^(l) from forward pass
    z_values: list of z^(l) from forward pass
    """
    m = 1  # single example
    L = len(weights)
    
    # Initialize gradients
    dW = [np.zeros_like(W) for W in weights]
    db = [np.zeros_like(b) for b in biases]
    
    # Output layer error
    delta = ((activations[-1] - y) * 
             sigmoid_derivative(activations[-1]))
    
    # Backpropagate
    for l in range(L-1, -1, -1):
        # Compute gradients for layer l
        dW[l] = np.outer(delta, activations[l])
        db[l] = delta
        
        # Propagate error to previous layer
        if l > 0:
            delta = (np.dot(weights[l].T, delta) * 
                    sigmoid_derivative(activations[l]))
    
    return dW, db

def train_network(X, y, weights, biases, learning_rate=0.01, 
                  epochs=1000):
    """
    X: training data (n_samples, 784)
    y: labels (n_samples, 10) - one-hot encoded
    """
    for epoch in range(epochs):
        total_cost = 0
        
        for i in range(len(X)):
            # Forward pass
            activations = [X[i]]
            z_values = []
            a = X[i]
            
            for W, b in zip(weights, biases):
                z = np.dot(W, a) + b
                z_values.append(z)
                a = sigmoid(z)
                activations.append(a)
            
            # Compute cost
            cost = np.sum((a - y[i])**2)
            total_cost += cost
            
            # Backpropagation
            dW, db = backpropagation(X[i], y[i], weights, 
                                     biases, activations, z_values)
            
            # Update weights and biases
            for l in range(len(weights)):
                weights[l] -= learning_rate * dW[l]
                biases[l] -= learning_rate * db[l]
        
        if epoch % 100 == 0:
            print(f"Epoch {epoch}, Avg Cost: {total_cost/len(X)}")
    
    return weights, biases
\end{lstlisting}

\subsection{Mini-Batch and Stochastic Gradient Descent}

\subsubsection{Full Batch vs Mini-Batch}

Full Batch:
\[
\nabla C = \frac{1}{n} \sum_{i=1}^{n} \nabla C_i \quad \text{[average gradient over ALL samples]}
\]
\begin{itemize}
    \item Expensive: $O(13,000 \times n)$ per step
    \item Accurate direction but slow
\end{itemize}

Mini-Batch (typical size: 32-256):
\[
\nabla C \approx \frac{1}{m} \sum_{i=1}^{m} \nabla C_i \quad \text{[average gradient over mini-batch]}
\]
\begin{itemize}
    \item Cheap: $O(13,000 \times m)$ per step (where $m \ll n$)
    \item Noisy gradient but fast convergence
    \item Good approximation in practice
\end{itemize}

Stochastic Gradient Descent (size 1):
\[
\nabla C \approx \nabla C_i \quad \text{[gradient of single example]}
\]
\begin{itemize}
    \item Very fast: $O(13,000)$ per step
    \item Noisy but allows online learning
\end{itemize}

\subsubsection{Algorithm}
\begin{lstlisting}
def train_sgd(X, y, weights, biases, learning_rate=0.01, 
              batch_size=32, epochs=10):
    """Stochastic Gradient Descent with mini-batches"""
    n = len(X)
    
    for epoch in range(epochs):
        # Shuffle data
        indices = np.random.permutation(n)
        X_shuffled = X[indices]
        y_shuffled = y[indices]
        
        # Process mini-batches
        for i in range(0, n, batch_size):
            X_batch = X_shuffled[i:i+batch_size]
            y_batch = y_shuffled[i:i+batch_size]
            
            # Accumulate gradients over mini-batch
            dW_batch = [np.zeros_like(W) for W in weights]
            db_batch = [np.zeros_like(b) for b in biases]
            
            for j in range(len(X_batch)):
                # Forward pass
                activations, z_values = forward_and_store(
                    X_batch[j], weights, biases
                )
                
                # Backpropagation
                dW, db = backpropagation(
                    X_batch[j], y_batch[j], weights, 
                    biases, activations, z_values
                )
                
                # Accumulate
                for l in range(len(weights)):
                    dW_batch[l] += dW[l]
                    db_batch[l] += db[l]
            
            # Average and update
            for l in range(len(weights)):
                weights[l] -= (learning_rate/len(X_batch)) * dW_batch[l]
                biases[l] -= (learning_rate/len(X_batch)) * db_batch[l]
    
    return weights, biases
\end{lstlisting}

\subsection{Why Backpropagation is Efficient}

\subsubsection{Naive Approach}
For 13,000 weights, computing each $\frac{\partial C}{\partial w}$:
\[
\frac{\partial C}{\partial w} = \frac{\partial C}{\partial a_{\text{output}}} \times \frac{\partial a}{\partial z} \times \cdots \times \frac{\partial z}{\partial w}
\]

Computing each derivative independently: $O(13,000)$

Computing all: $O(13,000^2) = O(169,000,000)$ per example!

\subsubsection{Backpropagation Solution}
Reuse intermediate computations:
\begin{enumerate}
    \item $\delta$ from one layer reused to compute $\delta$ for previous layer
    \item Share $\frac{\partial a}{\partial z}$ computations across multiple gradients
\end{enumerate}

Complexity: $O(13,000)$ per example!

Speedup: $13,000\times$

\newpage

\section{Summary and Key Formulas}

\subsection{Sigmoid Function}
\[
\sigma(z) = \frac{1}{1 + e^{-z}}
\]
\[
\sigma'(z) = \sigma(z)(1 - \sigma(z))
\]
Range: $(0, 1)$

\subsection{ReLU Function}
\[
f(z) = \max(0, z)
\]
\[
f'(z) = \begin{cases} 1 & \text{if } z > 0 \\ 0 & \text{otherwise} \end{cases}
\]
Range: $[0, \infty)$

Advantage: Easier to train, avoids vanishing gradient

\subsection{Cost Function (MSE)}
\[
C = \frac{1}{n} \sum_{i=1}^{n} \|\mathbf{y}_i - \mathbf{a}_i\|^2
\]

\subsection{Chain Rule (Backpropagation Core)}
\[
\frac{\partial C}{\partial w} = \frac{\partial C}{\partial a} \times \frac{\partial a}{\partial z} \times \frac{\partial z}{\partial w}
\]

\subsection{Full Forward Pass}
\[
\mathbf{a}^{(0)} = \mathbf{x} \quad \text{(input)}
\]
\[
\mathbf{z}^{(l)} = \mathbf{W}^{(l)} \mathbf{a}^{(l-1)} + \mathbf{b}^{(l)}
\]
\[
\mathbf{a}^{(l)} = \sigma(\mathbf{z}^{(l)})
\]
\[
\mathbf{y}_{\text{pred}} = \mathbf{a}^{(L)} \quad \text{(output)}
\]

\subsection{Full Backward Pass}
\[
\boldsymbol{\delta}^{(L)} = (\mathbf{y}_{\text{pred}} - \mathbf{y}_{\text{target}}) \odot \sigma'(\mathbf{z}^{(L)})
\]
\[
\boldsymbol{\delta}^{(l)} = (\mathbf{W}^{(l+1)T} \boldsymbol{\delta}^{(l+1)}) \odot \sigma'(\mathbf{z}^{(l)}) \quad \text{for } l = L-1, \ldots, 1
\]
\[
\frac{\partial C}{\partial \mathbf{W}^{(l)}} = \boldsymbol{\delta}^{(l)} (\mathbf{a}^{(l-1)})^T
\]
\[
\frac{\partial C}{\partial \mathbf{b}^{(l)}} = \boldsymbol{\delta}^{(l)}
\]

\subsection{Parameter Update}
\[
\mathbf{W}^{(l)} \leftarrow \mathbf{W}^{(l)} - \alpha \frac{\partial C}{\partial \mathbf{W}^{(l)}}
\]
\[
\mathbf{b}^{(l)} \leftarrow \mathbf{b}^{(l)} - \alpha \frac{\partial C}{\partial \mathbf{b}^{(l)}}
\]

Where $\alpha$ is the learning rate.

\newpage

\section{Practical Considerations}

\begin{enumerate}
    \item \textbf{Learning Rate}: Too high $\Rightarrow$ divergence; too low $\Rightarrow$ slow convergence
    \item \textbf{Weight Initialization}: Small random values ($\sim 0.01$) avoid saturation
    \item \textbf{Batch Size}: Trade-off between gradient accuracy and computational efficiency
    \item \textbf{Epochs}: Typically 10-100 for simple networks; more for complex data
    \item \textbf{Validation}: Always test on unseen data to check generalization
    \item \textbf{Regularization}: Add penalty terms (L1/L2) to prevent overfitting
    \item \textbf{Activation Functions}: Sigmoid (old), ReLU (modern), others for specific tasks
\end{enumerate}

\section{References}
\begin{itemize}
    \item Nielsen, M. A. (2015). Neural Networks and Deep Learning. \textit{Online resource}.
    \item Olah, C. (n.d.). Chris Olah's Blog. Retrieved from \url{http://colah.github.io/}
    \item Distill. (n.d.). Distill. Retrieved from \url{https://distill.pub/}
    \item Goodfellow, I., Bengio, Y., \& Courville, A. (2016). \textit{Deep Learning}. MIT Press.
\end{itemize}

\end{document}
